\section*{Bab \ref{ch:b3:origin-of-elements} --- Dari Mana Unsur-Unsur Berasal}
% Values from the Book 5 scratch calculator calc/ch30.py.

\begin{solution}{exo:b3:origin-of-elements:1}
(a)~$\mathrm{Si/Fe} = 10^{7.51 - 7.46} = 10^{0.05} = 1.12$ dan
$\mathrm{O/Fe} = 10^{8.69 - 7.46} = 10^{1.23} = 17.0$: sekitar satu atom silikon
dan tujuh belas atom oksigen per atom besi. (b)~$\mathrm{Fe/Au} =
10^{7.46 - 0.91} = 10^{6.55} = \num{3.55e6}$ dan $\mathrm{Pb/Au} =
10^{1.95 - 0.91} = 10^{1.04} = 11.0$: melewati puncak besi, kelimpahannya
turun hingga jutaan kali, tetapi timbal, yang berada di puncaknya sendiri, sebelas kali lebih melimpah
daripada emas.
\end{solution}

\begin{solution}{exo:b3:origin-of-elements:2}
(a)~$10^{7.46 - 12} = \num{2.884e-5}$, yaitu \num{28.8} atom besi per
juta atom hidrogen. (b)~Seratus kali lebih sedikit: \num{0.288} per
juta, satu atom besi per $1/\num{2.884e-7} = \num{3.47e6}$ atom
hidrogen. (c)~$[\mathrm{Fe/H}] = \log_{10}(\num{e-7}/\num{2.884e-5}) =
-2.46$.
\end{solution}

\begin{solution}{exo:b3:origin-of-elements:3}
$^{124}\mathrm{Xe}$ dan $^{126}\mathrm{Xe}$ terletak di sisi kaya proton,
di luar lintasan s, dan terlindung oleh $^{124}\mathrm{Te}$ dan
$^{126}\mathrm{Te}$ yang stabil: \omterm{def:b3:origin-of-elements:p-process}{inti p}. $^{128}\mathrm{Xe}$ dan $^{130}\mathrm{Xe}$
berada pada lintasan s, tetapi terlindung dari \omterm{def:b3:origin-of-elements:s-process}{proses r} oleh
$^{128}\mathrm{Te}$ dan $^{130}\mathrm{Te}$ yang stabil: hanya s. $^{129}\mathrm{Xe}$,
$^{131}\mathrm{Xe}$ dan $^{132}\mathrm{Xe}$ berada pada lintasan s dan tidak memiliki
isobar stabil dengan $Z$ lebih rendah (iodin hanya punya 127): kedua proses.
$^{134}\mathrm{Xe}$ dan $^{136}\mathrm{Xe}$ terletak di luar
$^{133}\mathrm{Xe}$, tempat lintasan s meninggalkan xenon, dan tidak
terlindung: hanya r.
\end{solution}

\begin{solution}{exo:b3:origin-of-elements:4}
(a)~$t = d/c = 40\times\qty{3.086e22}{m}/\qty{2.998e8}{m.s^{-1}} =
\qty{4.12e15}{s}$, sekitar \num{130} juta tahun. (b)~$\Delta v/c \le
1.7/\num{4.12e15} = \num{4.1e-16}$. Ini batas atas, karena sebagian
atau seluruh penundaan mungkin berasal dari pemancaran sinar gamma setelah
penggabungan.
\end{solution}

\begin{solution}{exo:b3:origin-of-elements:5}
(a)~$295\times0.1067 = \qty{31.48}{mb}$, $832\times0.0373 =
\qty{31.03}{mb}$, $155\times0.1984 = \qty{30.75}{mb}$. (b)~Dari
$^{122}\mathrm{Te}$: $N_{123} = 31.48/832 = \num{0.03783}$, \qty{1.4}{\percent}
di atas nilai terukur, dan $N_{124} = 31.48/155 = \num{0.2031}$,
\qty{2.4}{\percent} di atas. Kelimpahan yang berbeda lima kali lipat
dapat diprediksi hingga dalam beberapa persen hanya dari penampang lintang.
(c)~$^{128}\mathrm{Te}$ hanya dihasilkan oleh \omterm{def:b3:origin-of-elements:s-process}{proses r}: aliran s berhenti
pada tellurium di $^{127}\mathrm{Te}$, yang meluruh, sehingga tidak ada aliran s tunak yang
melewati $^{128}\mathrm{Te}$ dan kelimpahannya ditentukan oleh \omterm{def:b3:origin-of-elements:s-process}{proses r}.
\end{solution}

\begin{solution}{exo:b3:origin-of-elements:6}
(a)~$\tau_0\langle\sigma\rangle = 0.306\times155 = 47.4$, sehingga
$f = 1/(1 + 1/47.4) = \num{0.979}$ untuk $^{124}\mathrm{Te}$; untuk
$^{138}\mathrm{Ba}$, $0.306\times4.0 = 1.224$ dan $f = 1/(1 + 0.817) =
\num{0.550}$: inti ajaib hanya meneruskan \qty{55}{\percent} dari
aliran. (b)~Keenam faktor adalah \num{0.959}, \num{0.550}, \num{0.908},
\num{0.771}, \num{0.972} dan \num{0.915}; hasil kalinya \num{0.328}:
$\sigma N$ turun sebesar faktor \num{3.05}. (c)~Penurunan terukur adalah
$27.6/7.72 = 3.57$. Model memperkirakan besar loncatan dengan tepat dalam
\qty{15}{\percent}, hampir seluruhnya berasal dari $^{138}\mathrm{Ba}$ dan
$^{140}\mathrm{Ce}$, dua inti dengan $N = 82$.
\end{solution}

\begin{solution}{exo:b3:origin-of-elements:7}
(a)~$kT = 8\times10^3\times\qty{1.602e-19}{J} = \qty{1.282e-15}{J}$, dan
$v_T = \sqrt{2\times\num{1.282e-15}/\num{1.6749e-27}} =
\qty{1.24e6}{m.s^{-1}}$. (b)~Dengan $\langle\sigma\rangle =
\qty{155e-31}{m^2}$, $\lambda = n_\mathrm{n}\langle\sigma\rangle v_T$
adalah \qty{1.92e-11}{s^{-1}} untuk $n_\mathrm{n} = \qty{e12}{m^{-3}}$
($10^6$ per sentimeter kubik): waktu rata-rata \qty{5.2e10}{s}, sekitar
\num{1650} tahun; lalu \num{16.5} tahun pada $10^8$ dan \num{0.165}~tahun,
dua bulan, pada \qty{e10}{cm^{-3}}. (c)~$^{127}\mathrm{Te}$ meluruh dalam
hitungan jam, jauh lebih cepat daripada semua waktu tangkapan ini: ia selalu meluruh dan
lintasan tidak bercabang di sana. $^{134}\mathrm{Cs}$ hidup sekitar dua
tahun, berada di antara waktu tangkapan pada $10^8$ dan \qty{e10}{cm^{-3}}: pada
kerapatan tinggi dari kilatan helium, sebagian besar inti ini menangkap
neutron sebelum meluruh. Ia adalah titik percabangan, dan kelimpahan di kedua
sisinya mengukur kerapatan neutron.
\end{solution}

\begin{solution}{exo:b3:origin-of-elements:8}
(a)~$10^{6.30 - 3.14} = 1445$ kali lebih langka daripada kalsium dan
$10^{4.97 - 3.14} = 68$ kali lebih langka daripada titanium. (b)~$10^{8.46 - 1.38}
= \num{1.2e7}$. (c)~Skandium menunjukkan efek ganjil--genap: inti dengan
jumlah proton genap lebih kuat terikat dan dibangun oleh
penangkapan inti helium dalam tahap pembakaran bintang masif, yang
menguntungkannya. Berilium menunjukkan lembah litium, berilium dan boron:
inti-inti rapuh ini dihancurkan pada suhu interior bintang,
dan pembakaran helium langsung beralih dari helium ke karbon
tanpa menghasilkannya.
\end{solution}

\begin{solution}{exo:b3:origin-of-elements:9}
(a)~Misalkan $\alpha$ adalah \omterm{def:b3:origin-of-elements:alpha}{unsur alfa} dan $\mathrm{Fe}_\mathrm{cc}$,
$\mathrm{Fe}_\mathrm{Ia}$ besi dari kedua jenis supernova. Pada
dataran, $\alpha/\mathrm{Fe}_\mathrm{cc} = 10^{P}$ kali rasio Matahari.
Pada $[\mathrm{Fe/H}] = 0$, $\alpha/(\mathrm{Fe}_\mathrm{cc} +
\mathrm{Fe}_\mathrm{Ia})$ sama dengan rasio Matahari. Dengan membagi,
$\mathrm{Fe}_\mathrm{cc}/(\mathrm{Fe}_\mathrm{cc} + \mathrm{Fe}_\mathrm{Ia})
= 10^{-P}$, dan porsi supernova tipe~Ia adalah $1 - 10^{-P} = 1 - 10^{-0.4} =
\num{0.602}$. (b)~$10^{-P} = 1/3$ memberi $P = \log_{10}3 = 0.48$.
(c)~Supernova tipe~Ia mulai muncul setelah jeda yang sama di kedua galaksi.
Ketika mereka mulai di galaksi kerdil, gasnya baru mencapai
$[\mathrm{Fe/H}] = -2$, dibandingkan $-1$ di Galaksi: galaksi itu telah mengubah gasnya
menjadi bintang sepuluh kali lebih lambat.
\end{solution}

\begin{solution}{exo:b3:origin-of-elements:10}
(a)~Untuk zirkonium,
\[
\frac{(X/\mathrm{Mg})_\mathrm{Earth}}{(X/\mathrm{Mg})_\mathrm{CI}} =
\frac{10.5/228000}{3.96/95870} = 1.11 ;
\]
Demikian pula $(2670/228000)/(5010/95870) = 0.224$ untuk natrium,
\num{0.190} untuk kalium dan \num{0.118} untuk rubidium. (b)~Zirkonium, yang paling refraktori, berada
dekat dengan proporsi kondritnya. Ketiga logam alkali yang agak volatil,
yang mengembun pada \qtyrange{800}{1006}{K}, berkurang sebesar
faktor lima sampai sembilan, dengan rubidium yang paling volatil paling banyak. Bumi silikat
mengikuti \omterm{def:b3:origin-of-elements:volatility}{tren volatilitas}. (c)~Magnesium adalah unsur utama,
refraktori yang mengembun bersama sebagian besar batuan, dan ia bersifat
litofil: ia tetap berada di mantel ketika inti terbentuk. Unsur siderofil
sebagai acuan, seperti besi atau nikel, akan mencampur kehilangan ke inti dengan
kehilangan akibat volatilitas.
\end{solution}

\begin{solution}{exo:b3:origin-of-elements:11}
(a)~Ketika massa $\mathrm{d}s$ terkunci dalam bintang, angin menyapu
$\lambda\,\mathrm{d}s$ gas dengan \omterm{def:b3:origin-of-elements:metallicity}{metalisitas} $Z$:
$\mathrm{d}M_\mathrm{g} = -(1 + \lambda)\,\mathrm{d}s$ dan
$\mathrm{d}(ZM_\mathrm{g}) = y\,\mathrm{d}s - Z\,\mathrm{d}s -
\lambda Z\,\mathrm{d}s$. Dengan menguraikan persamaan kedua dan mengurangkan
$Z\,\mathrm{d}M_\mathrm{g}$ menyisakan $M_\mathrm{g}\,\mathrm{d}Z =
y\,\mathrm{d}s = -\frac{y}{1 + \lambda}\,\mathrm{d}M_\mathrm{g}$, yang
diintegralkan menjadi hasilnya. (b)~$\ln(1/0.29) = 1.238$: $Z = 0.0155$ untuk
$\lambda = 0$, dan setengahnya, $0.0077$, untuk $\lambda = 1$. (c)~Angin
menurunkan hasil efektif dengan faktor $1 + \lambda$. Dalam galaksi
bermassa kecil, supernova lebih mudah mendorong gas keluar, $\lambda$ besar,
dan logam ikut hilang bersama gas.
\end{solution}

\begin{solution}{exo:b3:origin-of-elements:12}
(a)~$N_\mathrm{Li}/N_\mathrm{H} = 10^{0.96 - 12} = \num{9.1e-12}$ di
Matahari dan $10^{3.25 - 12} = \num{1.78e-9}$ di kondrit CI. Matahari
menyimpan $10^{-2.29} = \num{0.0051}$ litiumnya: ia telah menghancurkan
\qty{99.5}{\percent} darinya. (b)~$\num{1.78e-9}/\num{1.45e-10} = 12.3$ dan
$\num{1.78e-9}/\num{4.72e-10} = 3.8$. (c)~Tata Surya menyimpan hampir
empat kali lebih banyak litium daripada yang dibuat Big Bang: sebagian besar litiumnya dibuat
belakangan, melalui proses lain. Bintang tertua menyimpan tiga kali
lebih sedikit daripada yang dibuat Big Bang: entah mereka telah menghancurkan sebagian darinya, seperti
Matahari, atau prediksinya melewatkan sesuatu; kedua perbandingan
saja tidak menentukan yang mana.
\end{solution}

\begin{solution}{pb:b3:origin-of-elements:1}

\textbf{Bagian I.}

\textbf{1.} $10^{12 - 7.46} = 10^{4.54} = \num{34700}$.

\textbf{2.} $N_\mathrm{He}/N_\mathrm{H} = 10^{10.914 - 12} = \num{0.0820}$, sehingga
$4N_\mathrm{He}/N_\mathrm{H} = \num{0.328}$, dibandingkan $Y/X =
0.2423/0.7438 = \num{0.326}$. Keduanya sepakat hingga \qty{0.7}{\percent}; selisih
kecil itu ada karena sebuah atom helium sedikit lebih ringan daripada
empat atom hidrogen.

\textbf{3.} $\log_{10}\bigl((1/350000)/\num{2.884e-5}\bigr) = -5.544 + 4.540 =
-1.00$.

\textbf{4.} $[\mathrm{O/H}] = -1.00 + 0.40 = -0.60$, sehingga
$\log\varepsilon(\mathrm{O}) = 8.69 - 0.60 = 8.09$. Bintang itu memiliki
$10^{8.09 - 6.46} = 42.7$ atom oksigen per atom besi, dibandingkan
$10^{1.23} = 17.0$ di Matahari: $10^{0.4} = 2.5$ kali lebih banyak.

\textbf{5.} $Z \approx 0.0139\times10^{-0.60} = \num{0.0035}$, seperempat
nilai Matahari.

\textbf{6.} Besi memiliki sangat banyak garis dalam spektrum bintang dingin, sedangkan
oksigen hanya memiliki sedikit dan samar: besi jauh lebih mudah
dan lebih teliti diukur, sehingga mewakili logam-logam lainnya.

\textbf{Bagian II.}

\textbf{7.} Masing-masing berada di belakang isobar stabil dengan $Z$ lebih rendah:
$^{122}\mathrm{Sn}$, $^{123}\mathrm{Sb}$, dan $^{124}\mathrm{Sn}$ menghentikan
rantai peluruhan \omterm{def:b3:origin-of-elements:s-process}{proses r} sebelum mencapai telurium. Ketiganya berada
pada jalur s, yang melewati telurium dari $^{122}\mathrm{Te}$ hingga
$^{126}\mathrm{Te}$.

\textbf{8.} \num{31.48}, \num{31.03}, dan \qty{30.75}{mb}; rata-rata
\qty{31.09}{mb}.

\textbf{9.} Rata-rata dari dua lainnya adalah $(31.48 + 30.75)/2 = \qty{31.11}{mb}$,
sehingga $N_{123} = 31.11/832 = \num{0.03740}$, hanya \qty{0.3}{\percent}
di atas nilai terukur \num{0.0373}.

\textbf{10.} $1.329\times44.4 = \qty{59.0}{mb}$, hampir dua kali lipat dataran.
$^{128}\mathrm{Te}$ hanya dibentuk \omterm{def:b3:origin-of-elements:s-process}{proses r}: jalur s meninggalkan telurium pada
$^{127}\mathrm{Te}$, sehingga kelimpahannya dibentuk oleh \omterm{def:b3:origin-of-elements:s-process}{proses r} dan
tidak berutang apa pun pada aliran s.

\textbf{11.} $1 - 7.43/8.53 = \num{0.13}$: \qty{13}{\percent} dari aliran melewati
$^{148}\mathrm{Sm}$, melalui tangkapan neutron pada
$^{147}\mathrm{Nd}$, $^{147}\mathrm{Pm}$, dan $^{148}\mathrm{Pm}$ yang tidak stabil sebelum
mereka dapat meluruh menjadi samarium.

\textbf{12.} Dengan kulit tertutup 82 neutron, inti-inti ini mengikat neutron tambahan
secara lemah dan memiliki penampang tangkapan yang sangat kecil (\qty{4.0}{mb}
untuk $^{138}\mathrm{Ba}$). Aliran menumpuk pada inti-inti ini dan hanya sebagian yang
diteruskan: inti ini adalah titik sempit.

\textbf{Bagian III.}

\textbf{13.} $\mu_1 = 13.7/(13.7 + 33.4) = 13.7/47.1 = \num{0.291}$.

\textbf{14.} $\ln(1/0.291) = \num{1.235}$, sehingga $y = 0.0154/1.235 = \num{0.0125}$,
\num{0.81} kali $Z_\odot$.

\textbf{15.} Pada $\mu = 0.9$, $\ln(1/0.9) = 0.1054$, sehingga $Z = 0.0125\times0.1054 =
\num{0.00131}$. Karena $0.00131/0.0154 = 0.0851$, $[\mathrm{Fe/H}] =
-1.07$. Pada $\mu = 0.5$, $Z = 0.0125\times0.693 = \num{0.00864}$ dan
$[\mathrm{Fe/H}] = -0.25$.

\textbf{16.} $Z = 0.00154$, sehingga $\mu = \mathrm{e}^{-0.00154/0.01247} =
\mathrm{e}^{-0.1235} = \num{0.884}$.

\textbf{17.} $1 - 0.884 = \num{0.116}$: \qty{12}{\percent} dari massa.

\textbf{18.} $y = 0.0154/\ln 10 = \num{0.0067}$, kira-kira separuhnya.

\textbf{Bagian IV.}

\textbf{19.} Massa bintang dengan \omterm{def:b3:origin-of-elements:metallicity}{metalisitas} di bawah $Z$ adalah
$M(1 - \mathrm{e}^{-Z/y})$, sehingga per satuan $Z$ nilainya
$(M/y)\,\mathrm{e}^{-Z/y}$. Karena $Z = Z_\odot10^{[\mathrm{Fe/H}]}$,
$\mathrm{d}Z = Z\ln 10\,\mathrm{d}[\mathrm{Fe/H}]$, dan per satuan
$[\mathrm{Fe/H}]$ distribusinya sebanding dengan
$Z\mathrm{e}^{-Z/y}$. Turunannya,
$(1 - Z/y)\mathrm{e}^{-Z/y}$, lenyap pada $Z = y$: puncaknya berada pada
$[\mathrm{Fe/H}] = \log_{10}0.810 = -0.09$.

\textbf{20.} $\num{0.116}/(1 - 0.291) = 0.1162/0.7091 = \num{0.164}$.

\textbf{21.} $Z/y = 1.235\times10^{-0.5} = 0.3905$, $1 - \mathrm{e}^{-0.3905} =
0.3233$, dan $0.3233/0.7091 = \num{0.456}$.

\textbf{22.} $6670\times0.1638 = 1093$, dibandingkan dengan 26: empat puluh kali terlalu banyak.

\textbf{23.} $26/902 = \qty{2.9}{\percent}$, tetap lebih dari lima kali di bawah
kotak. Bintang-bintang model adalah yang berumur panjang, yaitu yang massanya paling
besar sekitar satu massa Matahari, dan semuanya telah bertahan sejak terbentuk; di
antara bintang F yang lebih masif, yang tertua, yang juga paling miskin logam, sudah
mati, dan sampel lengkap kurang menghitung mereka.

\textbf{24.} $(1 - 0.1^{0.1})/(1 - 0.1) = 0.2057/0.9 = \num{0.23}$: lebih buruk.
Menaikkan fraksi gas pun tidak membantu: ketika $\mu_1$ mendekati 1,
fraksinya mendekati $0.1$, tetap lebih dari tiga kali
\qty{2.9}{\percent} yang teramati. Tidak ada fraksi gas yang dapat menyelamatkan kotak tertutup.

\textbf{25.} \omterm{thm:b3:origin-of-elements:closed-box}{Model kotak tertutup} memprediksi bahwa
\textbf{\qty[text-series-to-math=true]{16}{\percent}} bintang berumur panjang
di lingkungan Matahari memiliki $[\mathrm{Fe/H}] < -1$, dibandingkan dengan
paling banyak \qty{3}{\percent} yang teramati. Cakram tidak merupakan kotak tertutup:
ia tumbuh oleh aliran masuk gas miskin logam yang lambat, sehingga hampir tidak ada gas
yang pernah menunggu pada \omterm{def:b3:origin-of-elements:metallicity}{metalisitas} rendah untuk membentuk bintang.
\end{solution}
